\title{Byte Latent Transformer: Patches Scale Better Than Tokens}

\begin{abstract}
We present the Byte Latent Transformer ({\textsc BLT}), a novel architecture for large language models operating at the byte level. {\textsc BLT} achieves, for the first time, performance parity with traditional tokenization-based LLMs at scale, while substantially enhancing inference efficiency and model robustness. Rather than relying on fixed tokens, {\textsc BLT} groups bytes into variable-length patches, which form the basis for computation. The segmentation of these patches is guided by the entropy of the upcoming byte, thereby allocating greater computational resources and model depth to more complex, higher-entropy sequences. We provide the inaugural computationally-controlled scaling analysis of byte-level LLMs, scaling to 8B parameters and training on up to 4T bytes. Our findings confirm that scaling language models on raw bytes, absent a predetermined vocabulary, is a viable strategy. Efficiency in both training and inference is improved by adapting patch length in response to sequence predictability, leading also to qualitative gains in reasoning ability and outlier generalization. In summary, when inference cost is held constant, {\textsc BLT} exhibits considerably more effective scaling relative to models employing tokenization, benefiting from the joint increase of both patch size and parameter count.
\end{abstract}

\section{Introduction}
\label{section:intro}

We present the Byte Latent Transformer~(\textbf{BLT}), an architecture that forgoes traditional tokenization by directly modeling sequences of raw bytes. BLT is the first model of its kind to achieve parity, at scale, with standard tokenization-based approaches, while delivering notable gains in terms of both efficiency and robustness (\S\ref{sec:robustness}).
Most current large language models (llms) are trained using an almost end-to-end pipeline, with the exception of an initial tokenization stage—a non-learned, heuristic transformation that segments bytes into a predetermined vocabulary of tokens. This fixed tokenization injects biases into sequence compression, which introduces several limitations, including sensitivity to particular domains and modalities~\citep{dagangetting}, increased vulnerability to input noise~(\S\ref{sec:robustness}), deficits in orthographic generalization~\citep{edman2024cute}, and inequities across different languages~\citep{liang2023xlm,petrov2024language,limisiewicz2024myte}.

Historically, tokenization has played a crucial role since training large language models directly on raw bytes is extremely resource-intensive at large scales, largely because of the resulting extended sequence lengths~\citep{xue2022byt5}. To address this challenge, previous research has adopted either optimized self-attention strategies~\citep{el2020characterbert, clark2022canine} or utilized architectures that bypass attention entirely~\citep{wang2024mambabyte}~(\S\ref{section:related_work}). Nevertheless, these improvements predominantly benefit the training of \textit{smaller models}. For models at scale, the main computational burden in Transformers arises from the substantial feed-forward network layers that process every byte, rather than from the self-attention component.

To optimize computational resources, we introduce an adaptive, trainable technique for aggregating bytes into \textit{patches}~(\S\ref{section:patching}), along with a novel model architecture that integrates both byte-level and patch-level signals. In contrast to tokenization, {\textsc BLT} does not rely on a predetermined patch vocabulary. Instead, any selection of consecutive bytes can be projected into latent patch representations through compact encoder and decoder modules that are trained end-to-end. Our findings demonstrate that this approach yields a \textit{more} effective use of compute compared to models employing traditional tokenization schemes.

Tokenization-based LLMs uniformly distribute computational resources across all tokens, prioritizing uniformity over optimal efficiency. This approach relies on compression strategies that often do not align with the underlying complexity of token prediction tasks. Our model design instead emphasizes the importance of adaptively focusing computation where it is most necessary. For instance, predicting the final character in most words typically does not require the full power of a large transformer, as such tasks involve straightforward, low-entropy choices when compared to selecting the initial token of a new sentence. The architecture of {\textsc BLT} captures this principle~(\S\ref{section:architecture}), consisting of three transformer components: two compact byte-level \textit{local models} and one sizeable global \textit{latent transformer}~(\autoref{fig:arch_high_level}).

To facilitate dynamic computational allocation and to inform the grouping of bytes into patches, {\textsc BLT} organizes input data by segmenting it according to the entropy of subsequent byte predictions, thereby generating contextually sensitive clusters of bytes that exhibit approximately equal information density.

In this work, we conduct the first systematic scaling investigation of byte-level models controlled by flops, scaling up to 8B parameters and 4T training bytes. This demonstrates the feasibility of training large-scale models end-to-end from raw byte sequences without relying on fixed-vocabulary tokenizers. Notably, {\textsc BLT} achieves comparable performance to Llama 3 under matched training flop budgets\footnote{Model computational requirements are measured in Floating Point OPerations (flops).}, while reducing inference flops by as much as 50\%~(\S\ref{section:scaling}). 

Our analysis further reveals that operating directly on byte sequences markedly enhances performance on underrepresented, or long-tail, aspects of the data. {\textsc BLT} exhibits greater robustness to input noise when compared to models that use tokenizers, and shows superior character-level comprehension across orthographic reasoning, phonological patterns, and low-resource machine translation settings~(\S\ref{sec:robustness}).

Moreover, {\textsc BLT} enables joint scaling of both model size and patch size without surpassing the inference flop budget. By increasing patch length, the overall compute requirement decreases—since the transformer processes input less frequently—and this saved compute can be utilized to expand the global latent transformer. Through a series of inference-flop-controlled scaling experiments~(\autoref{fig:fixed_inference_scaling}), we find substantially improved scaling behavior relative to tokenization-based model architectures.

To conclude, this study presents the following key contributions: 1) We propose {\textsc BLT}, a byte-level latent LLM framework that adaptively distributes computational resources, thereby enhancing flop efficiency; 2) We demonstrate that our approach matches Llama 3 in training flop constraints up to the 8B parameter range, with the flexibility to exchange slight reductions in evaluation performance for up to 50\% improvements in flop efficiency; 3) {\textsc BLT} expands the scaling landscape for llms by enabling increases in model size without exceeding a predetermined inference compute ceiling; 4) We provide evidence that {\textsc BLT} exhibits greater resistance to input perturbations and possesses an understanding of sub-word input characteristics overlooked by traditional token-based llms. The implementation code for both training and inference of {\textsc BLT} is publicly available at~\url{https://github.com/facebookresearch/blt}.

\section{Patching: From Individual Bytes to Groups of Bytes}
\label{section:patching}

Dividing sequences of bytes into \textit{patches} enables {\textsc BLT} to adaptively assign computational resources according to contextual requirements. Various strategies for partitioning bytes into patches are illustrated in Figure~\ref{fig:patching_types}. More precisely, a patching function $f_p$ partitions a byte sequence $\pmb{x}=\{x_i,|i=1,\ldots n\}$ of length $n$ into $m$ patches ($m < n$), denoted as $\pmb{p}=\{p_j|j=1,\ldots,m\}$, by designating each $x_i$ as either 0 or 1, where a value of 1 signals the initiation of a new patch. In both token-centric and patch-centric architectures, the principal computational demand derives from the number of steps undertaken by the main Transformer block. For {\textsc BLT}, this corresponds to the total patches generated from the data given a specific patching approach. As a result, the mean length of a patch—referred to as the \textit{patch size}—is central to estimating the cost of model computation during training and evaluation for a particular patching function~(\S\ref{section:flops}). 
We next present three approaches to patching: fixed-length byte patching~(\S\ref{section:static-patch}), whitespace-based patching~(\S\ref{section:white-patch}), and adaptive patching informed by entropy metrics from a compact byte language model~(\S\ref{section:dyn-patch}). Lastly, we cover incremental patching techniques and clarify distinctions between tokenization and patching~(\S\ref{section:bpe-patch}).

\subsection{Strided Patching Every K Bytes}
\label{section:static-patch}

One common approach to byte grouping is to partition data into uniformly sized patches of $k$ bytes, as seen in MegaByte~\citep{yu2023megabyte}. Utilizing a constant stride simplifies both the implementation of training and inference and provides direct control over the average patch size, thus allowing easy management of the associated computational cost. Nonetheless, this method of patching introduces several notable limitations. Primarily, computation is not adaptively assigned based on local data complexity: computational resources may be squandered on uninformative regions (such as whitespace in code), or may be insufficient in areas that encode densely packed information, for example, in mathematical expressions. Furthermore, this approach often results in inconsistent and context-agnostic segmentation, such as identical words being split across patch boundaries in arbitrary or differing ways.

\subsection{Space Patching}
\label{section:white-patch}

\citet{slagle2024spacebyte} introduces an improvement to the strided patching method by generating new patches at the occurrence of space-like bytes\footnote{
    Space-like bytes are defined as any byte that is not a Latin character, a digit, or a UTF-8 continuation byte. Moreover, each patch must contain at least one non space-like byte. }, which frequently correspond to linguistically meaningful boundaries across various languages. In Space patching, a latent transformer step (entailing additional computation) is assigned to every word, guaranteeing consistent patching of words across different sequences while dedicating increased computational effort to the more challenging predictions that tend to appear immediately after spaces. For instance, in the prompt ``Who composed the Magic Flute?\uline{\hspace{.6em}}'', identifying the initial byte of the response is substantially more difficult than generating subsequent bytes once the answer begins (e.g., the ``M'' in ``Mozart''), as the first character sharply limits the pool of plausible continuations and makes the completion much more predictable. Nonetheless, Space patching faces limitations in applicability, especially across all languages and domains, and, crucially, lacks the ability to adjust patch size. To address these limitations, we next propose an alternative patching strategy based on the observation that predicting word-initial bytes is generally hardest, while simultaneously affording flexible control over the patch size.

\subsection{Entropy Patching: Using Next-Byte Entropies from a Small Byte LM}
\label{section:dyn-patch}

Instead of depending on rule-based heuristics like whitespace, we adopt a data-driven method to detect next-byte predictions with high uncertainty. We propose \textit{entropy patching}, a technique that leverages entropy estimates to determine the locations of patch boundaries.

We train a small byte-level auto-regressive language model on the training data for {\textsc BLT} and compute next byte entropies under the LM distribution $p_{e}$ over the byte vocabulary $\mathcal{V}$:

\begin{equation}
H(x_i) = \sum_{v \in \mathcal{V}} p_{e}(x_i=v|\pmb{x}_{<i}) \log p_{e}(x_i=v|\pmb{x}_{<i})
\end{equation}

We explore two approaches for detecting patch boundaries utilizing entropies $H(x_i)$. The initial method involves selecting points whose entropies exceed a predetermined global threshold, as shown in~\autoref{fig:patching}. Alternatively, the second strategy pinpoints points with entropy values significantly higher compared to the preceding point. This latter technique effectively detects locations that disrupt the otherwise approximately monotonically decreasing entropy within a patch.

\begin{align*}
    \text{Global Constraint}\hspace{1cm}H(x_t)& > \theta_g\\
    \text{Approx. Monotonic Constraint}\hspace{1cm}H(x_t)& - H(x_{t-1}) > \theta_r
\end{align*}

Patch boundaries are determined through a simple preprocessing procedure performed as part of the dataloading process. This contrasts with the approach of~\citet{nawrot-etal-2023-efficient}, in which a classifier is employed to infer patch boundaries based on entropy predictions. In our experimental setup~(\S\ref{section:experiments}), we evaluate and contrast both strategies for separating low- and high-entropy bytes.

\subsection{The Byte-Pair Encoding (BPE) Tokenizer and Incremental Patching}
\label{section:incremental-patching}
\label{section:bpe-patch}

A large number of contemporary llms, such as the baseline Llama 3 model, implement subword tokenization techniques like bpe~\citep{gage1994new, sennrich-etal-2016-neural}. In this context, ``tokens'' denote groups of bytes selected from a \textit{finite} set established before training, while ``patches'' pertain to adaptive groupings of byte sequences that are not constrained to any predefined vocabulary. The primary distinction between these concepts lies in the fact that, when using tokens, the model is unable to access the original byte-level information directly.

An important advancement offered by {\textsc BLT} compared to traditional tokenization-based methods is its novel approach to balancing vocabulary size and computational requirements. Conventional large language models (LLMs) typically face a situation where enlarging the vocabulary leads to, on average, bigger tokens and thus a reduced number of processing steps; however, this also results in a substantial increase in the size of the model's final projection layer. Consequently, tokenization-based systems are restricted in their ability to substantially alter token sizes or the computation needed at inference. For instance, Llama 3 manages to boost the average token length from 3.7 bytes to 4.4 bytes, but this comes at the expense of enlarging its embedding table by a factor of four when compared to Llama 2.

During generation, {\textsc BLT} must determine at each point in the byte sequence whether it is at the boundary of a patch; this choice is crucial because it dictates if additional computation using the Latent Transformer should be activated. Importantly, this decision must be made without any knowledge of the subsequent, as-yet-ungenerated portion of the sequence. Consequently, patch segmentation must rely solely on information from the past and current bytes, without reference to any future context. In formal terms, a patching method $f_p$ possesses the characteristic of incremental patching if:
$$f_p(\pmb{x}_{<i}) = f_p(\pmb{x})_{<i}$$
Notably, bpe does not qualify as an incremental patching scheme since its tokenization of an initial prefix can vary based on the tokens that follow, thus violating the condition specified above\footnote{Although introducing a dedicated delimiter token for marking patch boundaries could make bpe function incrementally, this approach leads to longer byte sequences.}.

\section{{\textsc BLT} Architecture}
\label{section:architecture}
The {\textsc BLT} system consists of a global, large-scale autoregressive language model that processes patch representations. This model is complemented by a pair of more compact local models: one responsible for converting byte sequences into patches, and another tasked with transforming patch representations back into bytes~(\autoref{fig:arch_high_level}).

\subsection{Latent Global Transformer Model}

\textit{The Latent Global Transformer} refers to an autoregressive transformer architecture $\mathcal{G}$ comprised of $l_{\mathcal{G}}$ layers. This model transforms an input sequence of latent patch embeddings $p_j$ into a corresponding sequence of output patch representations $o_j$. In this work, the notation $j$ designates individual patches while $i$ indexes bytes. The global transformer is equipped with a block-causal attention mechanism~\citep{dubey2024llama}, ensuring that attention for any given token is limited to its own and preceding patches within the document. As the most computationally demanding part of both pre-training and inference, the global model’s activation points can be modulated, allowing the allocation of computational resources across input and output regions to be tailored according to their respective complexities.

\subsection{Local Encoder}
\label{section:local}

\textit{The Local Encoder Model}, represented by $\mathcal{E}$, is a compact transformer-based model constructed with significantly fewer layers, $l_{\mathcal{E}} << l_{\mathcal{G}}$, compared to the global model. Its principal function is to rapidly convert a stream of input bytes $b_i$ into rich patch representations $p_j$. One notable modification to the standard transformer setup is the incorporation of a cross-attention layer succeeding each transformer layer, which aggregates byte-level representations into higher-level patch representations~(\autoref{fig:crossattn}).

Initially, each byte in the input sequence, $b_i$, is mapped via an embedding matrix $\mathbb{R}^{256 \times h_{\mathcal{E}}}$, resulting in embedded vectors $x_i$. These embeddings can be optionally enhanced with additional context by leveraging hash-embeddings, as described in~(\S\ref{sec:hash-emb}).

The model proceeds with a stack of alternating transformer and cross-attention layers~(\S\ref{sec:enc-cross-attn}), systematically refining the initial embeddings into patch representations, $p_i$, which subsequently serve as input for the global transformer, $\mathcal{G}$. The transformer component is equipped with a \textit{local block causal} attention mask, constraining each byte's attention to a predetermined window of $w_{\mathcal{E}}$ preceding bytes. While this window can span across dynamically defined patch boundaries, it remains restricted within the original document limits. Subsequent sections elaborate on the embedding strategy and the cross-attention mechanism.

\subsubsection{Encoder Hash n-gram Embeddings}
\label{sec:ngram-emb}
\label{sec:hash-emb}
An essential element for constructing effective and informative representations at every step $i$ involves embedding the context of prior bytes. Within {\textsc BLT}, this is accomplished by capturing the representation of the byte $b_i$ on its own \textit{as well as} in the context of an n-gram of bytes. For each position $i$, we begin by forming byte-grams

\begin{align}
    g_{i,n}=\{b_{i-n + 1},\ldots, b_i\}
\end{align}

for every byte position $i$, and for all $n$ ranging from three to eight.\footnote{
We disregard byte-grams of length $n$ or more in cases where $i<n$. }

Subsequently, we propose the use of hash $n$-gram embeddings, where all byte $n$-grams are projected via a hashing function to a particular index in a dedicated embedding table $E_{n}^{hash}$ of fixed dimension, for each $n$ in the set $\{3,4,5,6,7,8\}$~\citep{bai2010learning}.
The embedding obtained through this procedure is then summed to the corresponding byte embedding. This composite vector is normalized and subsequently supplied as input to the local encoder. The computation for the enriched embedding is defined by

\begin{align}
e_i &= x_i + \sum_{n = 3,...,8}E_{n}^{hash}(\text{Hash}(g_{i,n})) \\
\text{where, Hash}(g_{i,n}) &= \text{RollPolyHash}(g_{i,n}) \% |E_{n}^{hash}|
\end{align}

To compute $e_i$, we divide by the total number of considered $n$-gram sizes incremented by one, utilizing RollPolyHash as described in Appendix~\ref{appendix:rollpolyhash}. Section~\ref{section:ablations} presents an ablation study examining how varying $n$ and the size of the embedding table for $n$-gram hash embeddings impact trends in flop-constrained scaling laws. Beyond the use of hash-based $n$-gram embeddings, we also investigated frequency-based $n$-gram embeddings; further methodological details and findings related to this approach are discussed in Appendix~\ref{appendix:ngrams}.

\subsubsection{Encoder Multi-Headed Cross-Attention}
\label{sec:enc-cross-attn}
Our approach is largely based on the cross-attention input module from the Perceiver architecture~\citep{jaegle2021perceiver}; however, we diverge in a key aspect: whereas the original framework employs a static set of latent vectors, our method assigns latent representations dynamically to each patch (\autoref{fig:crossattn}), enabling them to attend solely to the constituent bytes of their respective patches. Within this module, each patch $p_j$ is represented by a unique query vector, constructed through pooling the byte embeddings for the relevant patch $p_j$, followed by application of a linear transformation, $\mathcal{E}_{C} \in \mathbb{R}^{h_{\mathcal{E}} \times (h_{\mathcal{E}} \times U_{\mathcal{E}})}$. Here, $U_{\mathcal{E}}$ represents the number of cross-attention heads in the encoder. To formalize this, let $f_{\text{bytes}}(p_j)$ denote the ordered bytes associated with patch $p_j$; then, we compute

\begin{align}
P_{0,j} &= \mathcal{E}_{C}(f_\text{bytes}((p_j)), f~ \text{is a pooling function} \\
P_l &= P_{l-1} + W_o\left(\text{softmax}\left(\frac{QK^T}{\sqrt{d_k}}\right)V\right) \\
\text{where } Q_j &= W_q(P_{l-1,j}), K_i = W_k(h_{l-1,i}), V_i = W_v(h_{l-1,i}) \\
h_l &= \text{Encoder-Transformer-Layer}_l(h_{l-1})
\end{align}

In this setup, $P \in \mathbb{R}^{n_p \times h_{\mathcal{G}}}$ denotes the collection of $n_p$ patch representations that serve as input to the global model; these are obtained by aggregating the byte embeddings $e_i$ associated with the constituent bytes of each patch $p_j$. The matrices $W_q$, $W_k$, $W_v$, and $W_o$ represent linear transformations used to generate the queries, keys, values, and output, respectively, with keys and values computed as projections of the byte-level representations $h_i$ from the preceding layer (or $e_i$ in the case of the initial layer). We apply a patch-specific masking approach whereby every query $Q_j$ is restricted to attend solely to the keys and values originating from the bytes in its corresponding patch $j$. Given that we perform multi-head attention over $Q, K$, and $V$—and since patch representations commonly occupy a higher-dimensional space ($h_{\mathcal{G}}$) in comparison to $h_{\mathcal{E}}$—we express $P_l$ as multiple attention heads, each with dimension $h_{\mathcal{E}}$, during the cross-attention computation, concatenating the resulting head outputs to recover the full $h_{\mathcal{G}}$ representation. Pre-LayerNorm is applied to the queries, keys, and values, and we do not incorporate any positional embeddings within this cross-attention mechanism. The module is further stabilized by adding a residual connection that encompasses the cross-attention layer.

\subsection{Local Decoder} 

In line with the structure of the local encoder, the local decoder $\mathcal{D}$ is also implemented as a compact transformer architecture, employing $l_{\mathcal{D}}$ layers, where $l_{\mathcal{D}} << l_{\mathcal{G}}$. Its primary function is to transform the global patch-level representations $o_j$ into the output sequence of raw bytes $y_i$. The local decoder generates each raw byte autoregressively, conditioned on the previously generated bytes, utilizing as input the hidden states from the local encoder corresponding to the current byte sequence. The architecture stacks $l_{\mathcal{D}}$ layers that alternate between cross-attention mechanisms and standard transformer operations. Within each layer, the cross-attention module precedes the transformer component, enabling the conversion of patch representations into detailed byte-level representations, after which the local decoder’s transformer sublayer processes the constructed byte sequence.

\subsubsection{Decoder Multi-headed Cross-Attention} 

In the decoder cross-attention, the roles of the queries and key/values are interchanged i.e. the byte-representations are now the queries, and the patch representations are now the key/values. The initial byte-representations for the cross-attention are initialized as the byte embeddings from the last encoder layer i.e. $h_{l_{\mathcal{E}}}$. The subsequent byte-representations for layer $l$, $d_{l,i}$ are computed as:

\begin{align}
D_0 &= h_{l_{\mathcal{E}}} \\
B_l &= D_{l-1} + W_o\left(\text{softmax}\left(\frac{QK^T}{\sqrt{d_k}}\right)V\right), \\
  &\text{where } Q_i = W_q(d_{l-1,i}), K_i = W_k(\mathcal{D}_C(o_j)), V_i = W_v(\mathcal{D}_C(o_j)) \\
  D_l &= \text{Decoder-Transformer-layer}_l(B_l)
\end{align}

In this formulation, the matrices $W_k$ and $W_v$ serve as projection matrices for keys and values, respectively, performing both linear transformation and a splitting operation $\mathcal{D}_C$ on the ultimate patch embeddings $o_j$ derived from the global model. The matrix $W_q$ acts as a query projection, applied to byte-level representations $d_{l-1}$ from the previous layer of the decoder transformer, or to $h_{l_{\mathcal{E}}}$ in the case of the initial layer. The output projection matrix is denoted by $W_o$. Consequently, $B$ is shaped as $B \in \mathbb{R}^{h_{\mathcal{D}} \times n_b}$, where $n_b$ designates the count of output bytes. The subsequent decoder representations, $D_l$, are obtained by passing the result of the cross-attention block $B$ through a layer of the decoder transformer. Consistent with the architecture used in the local encoder's cross-attention, we incorporate multiple attention heads, employ pre-LayerNorm, omit positional embeddings, and retain a residual connection that encompasses the cross-attention component.

\section{Experimental Setup}
\label{section:experiments}
We rigorously structure controlled experimental comparisons between {\textsc BLT} and conventional tokenization-based models, ensuring that {\textsc BLT} does not benefit from increased sequence context relative to its counterparts.

\subsection{Pre-training Datasets} In our experiments, models across all scales are initialized on two distinct datasets: 1) The Llama 2 dataset~\citep{touvron2023llama}, which consists of 2 trillion tokens sourced from a broad spectrum of publicly available corpora, then subjected to thorough cleaning and filtering procedures for quality enhancement; and 2) {\textsc BLT}-1T, a novel dataset comprising 1 trillion tokens curated from multiple public sources and further enriched with a selection of pre-training material from Datacomp-LM~\citep{li2024datacomp}. The Llama 2 dataset serves primarily in scaling law studies, leveraging the recommended token counts outlined by~\cite{dubey2024llama} to inform optimal design decisions for {\textsc BLT}, while {\textsc BLT}-1T is utilized for comprehensive pre-training to facilitate head-to-head comparisons with Llama 3 on evaluation benchmarks. Both datasets exclude content obtained from Meta platforms or services. For baseline studies involving tokenizer-based architectures, the Llama 3 tokenizer (vocabulary of 128K tokens) is employed, as it demonstrated superior baseline outcomes relative to the Llama 2 tokenizer in our tests.

\subsection{Entropy Model} For our experiments, the entropy model is implemented as a byte-level language model trained on the identical data distribution as the comprehensive {\textsc BLT} model. By default, this model utilizes a transformer architecture with 100M parameters, comprising 14 layers and a hidden size of 512, along with a sliding window attention mechanism covering 512 bytes. All other hyperparameters remain consistent with those used in our local and global transformer models. We have explored a variety of configurations, including alternate model sizes, receptive field lengths, and model architectures, as elaborated in \autoref{section:ablations}. Notably, when the model’s receptive field is sufficiently constrained, the entropy model can be stored efficiently as a lookup table.

\subsection{Entropy Threshold and Equalizing Context Length} When employing entropy-based patching methods, we determine an appropriate threshold to reach a specified average \textit{patch size} across the mixture of pretraining data. In {\textsc BLT}, unlike standard tokenization schemes, the \textit{patch size} is a flexible parameter, which directly affects the model's context length. To ensure parity in average context and prevent potential advantages from larger patch sizes, we fix the expected total byte count per batch. Accordingly, we shorten sequence lengths for models utilizing larger patch sizes. On the Llama 2 data, we adopt an 8k byte context, whereas for the {\textsc BLT}-1T dataset, the context is extended to 16k bytes on average, all while keeping the batch size constant at roughly 16M bytes.

Although the average batch size remains unchanged, the use of dynamic patching techniques produces varying proportions between bytes and patches in each batch. To improve efficiency, our {\textsc BLT} training implementation assembles batches so that they contain a fixed number of patches, thereby eliminating the need for padding within the computationally intensive latent transformer. Consequently, each batch consistently includes an identical count of patches. During model training, byte sequences are padded or truncated to 12k for Llama 2 and to 24k for {\textsc BLT}-1T datasets. This step mitigates memory spikes that could result from exceptionally large patches within sequences.

\subsection{Entropy Model Context}
\label{section:entropy-context}
Our observations indicate that applying entropy patching results in increasingly larger patches, especially within highly structured domains such as multiple choice tasks (as illustrated in \autoref{fig:mmlu_patching}), which tend to contain a high degree of repetition. These changes stem from reduced entropy in portions of the model context where content is repeated. To mitigate this, in our comprehensive evaluation of {\textsc BLT}-Entropy with a patch size of 4.5, we refresh the entropy context by inserting new lines and adopt an approximate monotonicity constraint. This constraint proves less vulnerable to issues of "entropy drift" that arise from variations in the length of the context. Notably, this adjustment only influences entropy calculation; the approach for determining the entropy threshold remains consistent.

\subsection{FLOPs Estimation} 
\label{section:flops}

Our methodology for calculating transformer flops is primarily based on the equations outlined in Chinchilla~\citep{hoffmann2022training}, which account for the flops associated with feed-forward networks, qkvo projections within self-attention, and both attention computation and output projection. However, in contrast to this prior work, we treat the input embedding layer as an efficient lookup rather than performing dense matrix multiplication, effectively considering it a 0-flop process. In accordance with established convention, we approximate the number of flops required for the backward pass to be double that of the forward pass.

To compute flops \textit{per byte} for {\textsc BLT} models, we add up the flops for the local encoder transformer, the global latent transformer, and the local decoder transformer, together with the cross attention blocks in the encoder and the decoder:

\begin{align}
    \text{FL}_{\text{{\textsc BLT}}} &= \frac{1}{n_p} \text{Transf. FL}(h_{\mathcal{G}}, l_{\mathcal{G}}, m=n_{ctx}/n_p, V=0) \\
   &\quad + \text{Transf. FL}(h_{\mathcal{E}}, l_{\mathcal{E}}, m=w_{\mathcal{E}}, V=0) \\
   &\quad + \text{Transf. FL}(h_{\mathcal{D}}, l_{\mathcal{D}}, m=w_{\mathcal{D}}, V=256) \\ 
    &\quad + \left[ \text{Cross Attn. FL}(h_{\mathcal{E}}, l_{\mathcal{E}}, m=n_p, r=n_p/k) \cdot \frac{k}{n_p} \right] \\
   &\quad + \text{Cross Attn. FL}(h_{\mathcal{D}}, l_{\mathcal{D}}, m=k, r=k/n_p) 
\end{align}

Here, $n_{ctx}$ denotes the sequence length measured in bytes, while $n_p$ refers to the patch size. The variable $r$ represents the proportion of queries to keys and values. The parameter $k$ specifies the ratio of the patch dimension to the byte dimension, indicating the number of local model components combined to create the full model representation (as illustrated in \autoref{fig:crossattn}, where $k=2$). $V$ stands for the size of the vocabulary used in the output projection, and this is relevant only in the context of the local decoder. The context length used in attention, $m$, varies depending on whether the operation is performed over bytes or patches. Accordingly, we adjust the calculation of attention FLOPs for each respective part. Detailed formulas for computing the FLOPs associated with Transformers and Cross-Attention are included in Appendix~\ref{section:flops-appendix}.

\subsection{Bits-Per-Byte Estimation} Perplexity only makes sense in the context of a fixed tokenizer as it is a measure of the uncertainty for each token. When comparing byte and token-level models, following previous work~\citep{xue2022byt5, yu2023megabyte, wang2024mambabyte}, we instead report Bits-Per-Byte (BPB), a tokenizer independent version of perplexity. Specifically:

\begin{align}
    \text{BPB}(x)&= \frac{\mathcal{L}_{CE}(\pmb{x})}{\ln(2)\cdot n_{\text{bytes}}}
\end{align}

Here, the aggregate cross-entropy loss quantifying the unpredictability in the data $\pmb{x}$ is divided by both the number of bytes comprising $\pmb{x}$ and a normalization constant.

\subsection{Transformer Architecture Hyperparameters} In configuring the transformer blocks within {\textsc BLT}—encompassing both its local and global variants—we primarily adopt the design choices implemented in Llama~3~\citep{dubey2024llama}. Specifically, all feed-forward modules utilize the SwiGLU activation function~\citep{swiglu}, while rotary positional embeddings (RoPE)~\citep{RoPe} are applied exclusively within the self-attention layers, with $\theta=500000$ as described in \citet{xiong2024effective}. For layer normalization, RMSNorm~\citep{RMSNorm} is employed across the model. The self-attention mechanisms rely on Flash attention~\citep{dao2022flashattention} when operating with attention masks that are either block causal or adhere to a fixed-window block causal structure; the fixed-width attention configurations employ a window size of 512. In contrast, for cross-attention modules—where masks are dynamically determined based on input patches—we deploy Flex Attention\footnote{\url{https://pytorch.org/blog/flexattention}} to facilitate efficient, fused operations, which greatly accelerates the training process.

\subsection{{\textsc BLT}-Specific Hyperparameters} In order to assess the performance of {\textsc BLT} models, our experiments span two main approaches: analyzing scaling behavior and evaluating performance on downstream tasks. These studies utilize models of varying capacities, specifically at the 400M, 1B, 2B, 4B, and 8B parameter scales. Detailed configuration parameters for each model architecture can be found in Appendix~Table~\ref{tab:arch}.  
For initialization of the queries in the first cross-attention layer within the local encoder, we employ max-pooling. Consistently across all {\textsc BLT} models, we implement $500,000$ hashes using one hash function, and utilize n-gram sizes between 3 and 8. All models are trained with a learning rate of $4e-4$. Selection of this shared learning rate for both token and {\textsc BLT} model variants is guided by a hyperparameter search conducted over the range $1e-3$ to $1e-4$ at model sizes 400M and 1B, indicating that this learning rate yields optimal results. For investigating scaling behavior using Llama-2 data, we set the training batch sizes to follow the guidance in~\cite{dubey2024llama}, or use an equivalent batch size measured in bytes.  
We utilize the AdamW optimizer~\citep{adamw}, setting $\beta_1$ to 0.9, $\beta_2$ to 0.95, and $\epsilon=10^{-8}$. The learning rate schedule incorporates a linear warm-up phase of 2000 steps, followed by a cosine decay to zero. Additionally, a weight decay factor of 0.1 is applied, together with global gradient clipping at a value of 1.0.

\section{Scaling Trends}
\label{section:scaling}

We provide a comprehensive analysis of the scaling patterns exhibited by byte-level models, with the goal of offering insights relevant to the future scaling of {\textsc BLT} architectures. Our scaling analysis seeks to overcome several shortcomings identified in earlier investigations of byte-level models by: (a) systematically examining trends within the compute-optimal training paradigm, (b) developing parallel 8B parameter models trained on substantial datasets—reaching as much as 1T tokens or 4T bytes—and subsequently assessing their performance across various downstream tasks, and (c) evaluating scaling behavior under conditions where inference cost is held constant. In subsequent sections, we will explore particular benefits that arise from modeling data at the byte-sequence level.

\subsection{Parameter Matched Compute Optimal Scaling Trends}
\label{section:parameter-matched-scaling}
Employing the Llama 2 dataset, we conduct training runs of several \textit{compute-optimal} bpe and {\textsc BLT} models at four distinct parameter scales, spanning from 1B to 8B. Subsequently, we chart the relationship between total training flops and language modeling accuracy using a representative portion of the composite training dataset. The bpe architectures are configured and trained with the empirically optimal proportion between model parameters and the quantity of training data, following guidance from Llama~3~\citep{dubey2024llama}. This methodology is theoretically expected to maximize performance given fixed training compute resources~\citep{hoffmann2022training}, thus providing a robust benchmark. For each bpe configuration, an equivalent {\textsc BLT} model—incorporating a Latent Transformer of identical parameter count and architecture—is trained using the exact same dataset.

As shown in~\autoref{fig:scaling-compute-opt} (right), {\textsc BLT} models consistently equal or surpass their bpe equivalents, and this performance pattern persists across increasing model sizes and computational resources. To our knowledge, {\textsc BLT} represents the first byte-level Transformer architecture to attain comparable scaling behavior with BPE-based models in compute-optimal settings. This observation supports our hypothesis that the optimal parameter-to-compute ratio established for bpe is applicable to {\textsc BLT} as well, or at the very least, closely approximates it.

Both advances in architecture and the introduction of dynamic patching are essential for aligning with bpe scaling behavior. In ~\autoref{fig:scaling-compute-opt} (left), we present a comparison of models utilizing space patching with Llama 3. SpaceByte \citep{slagle2024spacebyte} is estimated by employing {\textsc BLT} space patching, but without incorporating n-gram embeddings and cross-attention mechanisms. While SpaceByte demonstrates improvement over Megabyte, it does not reach the performance level of Llama 3. As shown in \autoref{fig:scaling-compute-opt} (right), gains are achieved when integrating both new architectural components and dynamic patching strategies. When sufficiently scaled, {\textsc BLT} models match the performance of leading tokenizer-based systems like Llama 3.

Additionally, we find that for tokenizer-based models, the selection of tokenizer plays a significant role in model effectiveness; specifically, models utilizing the Llama-3 tokenizer achieve higher performance than those employing the Llama-2 tokenizer when both are trained on identical data.

In summary, the {\textsc BLT} architecture positions itself between Llama 2 and Llama 3, particularly when operating with much larger patch sizes. The mean token length for the bpe tokenizers in Llama 2 and 3 is 3.7 and 4.4 bytes, respectively. By comparison, {\textsc BLT} attains comparable scaling behaviors when utilizing average patch sizes of 6 or even 8 bytes. Since inference flop counts decrease as average patch size increases, adopting an 8-byte patch size can yield approximately 50\% reduction in inference flop. Moreover, models using larger patch sizes tend to exhibit improved performance as both the model and data scale up. For instance, {\textsc BLT} with an 8-byte patch size initially underperforms relative to the bpe Llama 2 at the 1B scale but surpasses bpe performance at 7B parameters. These observations point toward the potential benefits of employing even larger patch sizes at higher model and data scales, and imply that increasing patch sizes further may be viable as computational resources expand.

\subsection{Beyond Compute Optimal Task Evaluations}
\label{section:evals}
To further investigate scaling characteristics, we extend the training of an 8B {\textsc BLT} model past the compute-optimal point, using the more extensive and higher-quality {\textsc BLT}-1T dataset. Performance is assessed through a collection of widely-used classification and generative benchmarks. For evaluation purposes, we focus on tasks spanning common sense reasoning, world knowledge, and code generation:
\paragraph{Classification tasks} encompass ARC-Easy (0-shot)~\citep{Eval_ARC}, Arc-Challenge (0-shot)~\citep{Eval_ARC}, HellaSwag (0-shot)~\citep{Eval_hellaswag}, PIQA (0-shot)~\citep{Eval_piqa}, and MMLU (5-shot)~\citep{hendrycks2020measuring}. For these, we utilize a prompt-scoring approach that determines likelihood over possible answer choices, ultimately reporting mean accuracy across tasks.
\paragraph{Coding related generation tasks:} We present pass@1 metrics for both MBPP (3-shot)~\citep{Eval_MBPP} and HumanEval (0-shot)~\citep{Eval_HumanEval} to gauge the proficiency of LLMs in producing correct Python code.

In \autoref{tab:evals}, we present a comparative analysis of three models, each trained using the {\textsc BLT}-1T dataset: a bpe Llama 3 tokenizer-based model\footnote{We selected the Llama 3 tokenizer, featuring a 128k vocabulary, due to its superior performance relative to the 32k vocabulary of Llama 2.}, along with two versions of the {\textsc BLT} model—one integrating a space-patching approach ({\textsc BLT}-Space), and the other leveraging an entropy-based patching approach ({\textsc BLT}-Entropy). For the entropy model, we apply an approximate monotonicity constraint and refresh the model’s context at each new line, as outlined in~\autoref{section:entropy-context}. All three configurations are trained using comparable flop budgets. Notably, for {\textsc BLT}-Entropy, we further adjust the entropy threshold at inference, reducing it from 0.6 to 0.1—a modification that enhances task accuracy but increases the number of inference steps.

The {\textsc BLT}-Entropy model achieves superior results compared to the Llama 3 model on 4 out of 7 benchmarks, despite both being trained with an equivalent number of bytes. This performance advantage likely arises from two factors: (1) more efficient utilization of training computation enabled by dynamic patching, and (2) the explicit modeling of information at the byte level rather than at the token level.

Conversely, {\textsc BLT}-Space lags behind the Llama 3 tokenizer on all tasks except one, although it notably lowers inference flops due to its larger mean patch size of 6 bytes. In contrast, both the bpe and entropy-patching models utilize an average patch size near 4.5 bytes across the training data mix. Given an identical training budget, the model with a greater patch size is able to process 30\% more data than the other two approaches, which may consequently move {\textsc BLT} further from achieving compute-optimality.

\subsection{Patches Scale Better Than Tokens} In the context of {\textsc BLT} models, it is feasible to expand both the model size and the patch size concurrently without surpassing the allocated flops for training and inference, nor increasing the volume of training data. Patch-based approaches, distinct from models that use a fixed vocabulary of tokens, allow for arbitrary scaling of patch size. This flexibility overcomes the efficiency constraints characteristic of token-based systems, as elaborated in Section~\ref{section:bpe-patch}. Using larger patch sizes leads to computational savings, which can then be redirected to enhance the capacity of the global latent transformer since it will require fewer executions.

We perform a scaling experiment at a constant inference cost to examine whether increasing model size while reducing the number of steps with larger patches yields better performance compared to employing smaller models that require more steps. Using the Llama 2 tokenizer, we begin with models of 400 million and 3.6 billion parameters, then identify flop-equivalent models that utilize the Llama 3 tokenizer as well as {\textsc BLT}-Entropy models with mean patch sizes of 6 and 8 bytes based on the training datamix (see~\autoref{tab:fixed-inf-params} for specific model configurations). For models utilizing an 8-byte patch size, we opt to implement 3 encoder layers instead of just 1. Each model is trained across a range of allocated training flop budgets.

As depicted in \autoref{fig:fixed_inference_scaling}, {\textsc BLT} models display superior scaling behavior compared to tokenization-based counterparts across both inference flop categories. Notably, while BPE models tend to excel when training budgets are limited, they are swiftly outperformed by {\textsc BLT} models just past the compute-optimal point. In practical terms, it can often be advantageous to invest additional resources during initial pretraining to obtain a higher-performing model for a given inference cost. The 8B model class provides a clear illustration of this trend; for instance, models like Llama 3.1 have been pretrained on data volumes two orders of magnitude greater than the compute-optimal quantity for their parameter count.

The threshold at which {\textsc BLT} surpasses token-based architectures has moved somewhat nearer to the compute-optimal regime as the model size increases, shifting from 3x to approximately 2.5x the optimal compute budget in the higher flop category. Additionally, the model utilizing patch size 8 displays a more pronounced scaling behavior among the larger flop models, allowing it to outperform other model configurations more rapidly. As outlined in Section~\ref{section:parameter-matched-scaling}, using larger patch sizes tends to bring performance metrics closer to those observed in BPE-based models as the model scale increases. We partially attribute this pattern to the diminishing proportion of total computational cost incurred by the byte-level Encoder and Decoder modules, which do not scale as quickly as the Latent Transformer. When expanding the total parameter count by a factor of 20 (from 400M to 8B), the local parameters specific to {\textsc BLT} increase by roughly twofold. This distinction is significant because enlarging the patch size mainly impacts the computational load of the patch Latent Transformer, leaving the flops associated with the byte-level modules largely unchanged. Consequently, the {\textsc BLT}-Entropy model with patch size 8 saw its size rise from 1.6x to 1.7x that of the Llama 2 model when scaled up to the larger parameter setting.

To summarize, our investigation into scaling patch lengths indicates that the {\textsc BLT} patch-based architecture exhibits improved scaling behavior when patch size and model capacity are augmented together. This advantageous trend persists and, in fact, becomes more pronounced as the models increase in size.

\section{Byte Modeling Improves Robustness} Additionally, we assess the resilience of {\textsc BLT} relative to models utilizing tokenization without direct byte-level input, and we introduce a method to enable pretrained token-based models to process byte representations.  
\label{sec:robustness}

\subsection{Character-Level Tasks}

One of the initial reasons for exploring byte-level models was their inherent resilience to noise at the byte level in the input, as well as their capacity to recognize the internal structure of tokens—an area where models relying on tokenizers typically fall short. To assess these qualities, we conduct supplementary experiments using benchmarks designed to test both resistance to input perturbations and sensitivity to character-level features, encompassing not only English and multiple languages but also numeric digits and phonemes. The outcomes of these evaluations are reported in Table~\ref{tab:char_tasks}.

\paragraph{Noisy Data} To assess and contrast the robustness of tokenizer-based models with {\textsc BLT}, we generate noised variants of the standard classification benchmarks outlined in Section~\ref{section:evals}. Five unique character-level noising techniques are utilized to alter the original text: (a) \textit{AntSpeak}, where text is rendered entirely in uppercase with spaces between each character; (b) \textit{Drop}, which deletes 10\% of the characters at random; (c) \textit{RandomCase}, applying random capitalization such that half of the characters become uppercase and the remainder lowercase; (d) \textit{Repeat}, where 20\% of the characters are duplicated up to four times; and (e) \textit{UpperCase}, which capitalizes all characters in the input. Each noising method is applied either to the prompt, the completion, or both, considered separately during evaluation, and average results are reported. Table \ref{tab:char_tasks} presents findings from noisy HellaSwag~\citep{Eval_hellaswag}, showing that {\textsc BLT} consistently achieves superior robustness compared to tokenizer-based models. Specifically, {\textsc BLT} exceeds the performance of models trained on identical datasets by an average margin of 8 points, and also surpasses the Llama 3.1 model, even though Llama was trained on a substantially larger corpus.

\paragraph{Phonology - Grapheme-to-Phoneme (G2P)} We evaluate {\textsc BLT}'s performance in converting sequences of graphemes (individual characters making up a word) into corresponding phonemic representations, which capture the word's pronunciation. The results for the G2P task in a 5-shot scenario, obtained with Phonology Bench~\citep{suvarna2024phonologybench}, are shown in Table \ref{tab:char_tasks}. These findings indicate that {\textsc BLT} exceeds the performance of the Llama 3 1T tokenizer-based baseline model on this specific task.

\paragraph{CUTE}
To evaluate character-level comprehension, we test {\textsc BLT} using the CUTE benchmark~\citep{edman2024cute}. This benchmark consists of various tasks that are generally divided into three main types: compositional understanding, orthographic similarity, and sequence manipulation capabilities. Most tokenizer-based models find this benchmark particularly difficult, as they seem to be aware of the spellings of their individual tokens but lack the capacity to leverage this knowledge effectively for text manipulation. According to Table~\ref{tab:char_tasks}, {\textsc BLT}-Entropy surpasses both BPE Llama 3 models by over 25 points on these evaluations. Notably, our model excels at tasks involving character manipulation, attaining a score of 99.9\% on the spelling-focused assessments. The magnitude of these gains, especially considering {\textsc BLT} was trained with only 1/16th the data compared to Llama 3.1, suggests that learning character-level features is a substantial challenge for BPE-based approaches. Figure \ref{fig:cute_examples} provides representative examples in which the Llama 3 tokenizer model falters, whereas the {\textsc BLT} model succeeds. The only tasks where BPE models outperform are word deletion and insertion. While byte-level models may not handle such word-level changes as directly, the performance gap remains small, implying that constructing words from characters could prove less difficult than the reverse. For all tasks, we retain a consistent evaluation protocol using the original prompts from Huggingface. Additionally, BPE models might show improved results with further prompt engineering.

\paragraph{Low Resource Machine Translation} We assess the performance of {\textsc BLT} on both directions of translation across six major language families and twenty-one low-resource languages, each featuring a range of scripts from the FLORES-101 benchmark~\citep{goyal2022flores}. SentencePiece BLEU scores are presented in Table~\ref{tab:flores}. The evaluation reveals that {\textsc BLT} consistently surpasses the model utilizing the Llama 3 tokenizer, registering an average improvement of 2 BLEU points for translations into English and 0.5 BLEU points for translations from English. For more commonly studied language pairs, {\textsc BLT}'s performance is comparable to or marginally better than that of Llama 3. Notably, the gains are more pronounced for low-resource language pairs, where {\textsc BLT} outperforms Llama 3 by a wider margin, highlighting the benefits of byte-level modeling in effectively handling diverse and infrequent byte sequences.

\subsection{Training {\textsc BLT} from Llama 3}
\label{sec:distillation}
We investigate a procedure in which {\textsc BLT} models benefit from existing pre-trained models based on tokenization by initializing {\textsc BLT}'s global transformer parameters with those from a pre-trained Llama 3.1 checkpoint. In this setup, the global transformer weights are further trained with a learning rate set to one-tenth of that used for the local encoder and decoder, applying the same total training steps as used for Llama 3. Performance outcomes in Table~\ref{tab:distill} demonstrate that {\textsc BLT} initialized from Llama 3.1 yields results substantially better than both the Llama 3 and baseline {\textsc BLT} models, each of which was trained under the same computational budget (in FLOPs). Furthermore, comparing these results with our {\textsc BLT}-Entropy configuration (from Table~\ref{tab:evals}), which was exposed to a much larger corpus of 1T tokens, the {\textsc BLT} initialized from Llama 3.1 continues to achieve higher scores on the MMLU benchmark. This suggests the method offers an efficient alternative to drastically lower required training FLOPs.

This approach can alternatively be interpreted as adapting tokenizer-based architectures to operate without a tokenizer, thereby reconfiguring a pre-trained LLaMA 3.1 model into a {\textsc BLT} model. For a thorough evaluation, we present both the baseline LLaMA 3.1, originally trained with 15T tokens, and the {\textsc BLT} variant derived from LLaMA 3 in Table \ref{tab:distill}. Our results indicate that while the new model suffers only a minor reduction in performance on MMLU and HumanEval, there is a more pronounced decrease across additional benchmarks. These findings highlight the necessity of continued research to more effectively exploit the capacity of the pre-trained model, particularly through refinement of data compositions and hyperparameter tuning.

\section{Ablations and Discussion}
\label{section:ablations}
Here, we present ablation studies that validate the architectural decisions underlying {\textsc BLT}, as well as the selected patching strategy and hyperparameter configurations for the {\textsc BLT} model with 8B parameters trained on the {\textsc BLT}-1T dataset.

\paragraph{Entropy Model Hyper-parameters} In order to analyze how the size of the entropy model and the length of the context window impact scaling, we conduct experiments training byte-level entropy transformer models that range from 1m to 100m parameters, using context windows that span from 64 to 512 tokens. We then generate plots depicting bpb versus training flop scaling law curves, based on our $400m$ and $1b$ {\textsc BLT} models that have been trained on the Llama-2 dataset, and display these results in \autoref{fig:entropy_ablation}. The results show a positive correlation between scaling performance and increases along both the model size and context window dimensions, although gains plateau when the model exceeds 50m parameters.

\paragraph{Types of Patching}

We conduct an ablation study of the four patching strategies outlined in Section~\ref{section:patching}, which include: (1) Strided Patching utilizing strides of 4 and 6, (2) Whitespace-based Patching, (3) Patching according to BPE Tokenizer segments leveraging the Llama 3 tokenizer, and (4) Entropy-driven patching employing a compact byte-level language model.

Although dynamic patching shortens the actual length of processed sequences, we standardize sequence length across all patching strategies to ensure comparable context sizes. This approach guarantees that, on average, every model encounters the same number of bytes per sequence throughout both training and inference phases, thereby eliminating potential confounds related to varying context modeling capacities. Figure~\ref{fig:scaling-compute-opt} presents the outcomes from these controlled experiments. Among all configurations, every patching approach except static patching achieves superior performance, with space patching showing nearly equivalent results to dynamic entropy-based patching.

In \autoref{tab:evals_ablation}, we report on benchmark evaluations for {\textsc BLT} models, contrasting approaches based on traditional tokenizers, space patching, and entropy-based patching. All models were trained on the Llama 2 dataset for the number of steps identified as optimal in~\citep{dubey2024llama}. While space patching constitutes a comparatively straightforward method—bypassing the need to invoke an entropy model dynamically during training—we observe that the benefits associated with entropy-based patching in our scaling experiments (see Section~\ref{section:scaling}) persist and are observable in the context of downstream benchmark tasks as well.\footnote{The results for space patching correspond to prior experiments lacking cross-attention, but the main qualitative patterns remain unchanged when cross-attention is included.}

\paragraph{Cross-Attention} 
In~\autoref{tab:crossattn_ablations_1b}, we examine the effects of introducing cross-attention at different locations within both the encoder and decoder components of {\textsc BLT}. Regarding encoder cross-attention, we explore several methods for initializing the queries: (1) using an identical learned embedding across all global states, (2) employing a hash embedding derived from the bytes present in the patch, and (3) utilizing a pooled summary of the encoder’s hidden representations for the patch bytes at the respective encoder layer.

Our results indicate that implementing cross-attention within the \textit{decoder} yields the greatest benefit. Within the encoder, cross-attention brings a modest gain, but this advantage is observed solely when queries are initialized through pooling. Furthermore, our analysis shows that cross-attention is particularly advantageous for the Common-Crawl dataset, with the effect becoming more pronounced as patch sizes increase.

\paragraph{n-gram Hash Embeddings} 

We evaluate the impact of varying n-gram hash embedding vocabulary sizes (0, 100K, 200K, and 400K) and report our findings in Table~\ref{tab:ngram_ablations_1b}. Our analysis demonstrates that hash embeddings yield improvements across all evaluated domains, with notably greater enhancements observed for Wikipedia and Github datasets (a 0.04 bpb gain versus a 0.01 bpb gain after 15k steps at the 8B scale). Reducing the number of hashes from 500K to 300K at the 8B model size results in a negligible performance decrease of approximately 0.001 bpb after 15k steps. These results suggest that the use of hashes is crucial for achieving parity between {\textsc BLT} and tokenizer-based model performance; nonetheless, the incremental benefits decline when surpassing 300K hashes. Moreover, the observed improvements from hashing and cross-attention mechanisms are largely additive, as they tend to contribute distinct benefits across different datasets.

\paragraph{Local Model Hyperparamaters} Table \ref{tab:local_layers} presents an ablation of different configurations concerning the layer counts in the local encoder and decoder. When utilizing hash n-gram embeddings, {\textsc BLT} demonstrates strong performance with a highly minimal encoder—specifically, one consisting of a single layer—while benefitting from a more complex decoder.

\section{Related Work}
\label{section:related_work}

\textbf{Character-Level RNNs:} The task of Character Language Modeling has maintained its relevance from the outset of neural approaches~\citep{sutskever2011generating,mikolov2012subword,graves2013generating}, primarily because such models naturally accommodate out-of-vocabulary terms without requiring fallback strategies. \cite{kim2016character} introduce a framework where input sequences are composed solely of characters, utilizing convolutional and highway architectures to generate representations passed into LSTM-based RNNs; their results match then-current leading RNN language models for English and demonstrate superior performance for languages with rich morphological structures—a long-standing motivation for character-level LLMs. Similarly, \cite{kenter2018byte} leverage byte-level LSTM architectures for machine comprehension, which consistently outperform word-level baselines, especially for highly inflected languages like Turkish and Russian. In a related approach, \cite{zhang2015character} employ character-centric convolutional networks for classification, observing that these models surpass word-based alternatives in certain scenarios. Expanding upon such models, \citet{chung2019hierarchical} implement hierarchical LSTMs, wherein boundary detectors at each hierarchy identify latent textual structure, achieving further gains on character-level language modeling tasks. By contrast, ByteNet~\cite{kalchbrenner2016neural} adopts a CNN-based paradigm applied to character sequences, offering an alternative to attention mechanisms in machine translation applications.

\textbf{Character-Level Transformers:} The introduction of attention mechanisms in transformer architectures~\citep{vaswani2017attention} alongside subword tokenization methods~\citep{sennrich-etal-2016-neural} has led to notable advancements in the capabilities of neural models for language modeling and standardized evaluation tasks. Nevertheless, employing word or subword segments inherently imposes an inductive bias regarding the granularity at which these models process linguistic data. Seeking to bridge the advantages provided by transformer-based approaches with the promising early outcomes demonstrated by character-level language models, \citet{al2019character} leveraged very deep transformer networks supplemented by auxiliary loss functions. Through this approach, they were able to develop transformer language models at the character level that surpassed earlier LSTM-based counterparts in performance. Despite these gains, a substantial disparity relative to word-level language models persisted. In line with these observations, GPT-2~\citep{radfordgpt2} found that for expansive datasets such as the 1 billion word benchmark, byte-level language models failed to match the effectiveness of those trained on word-level units.

Although \cite{Choe2019BridgingTG} showed that transformer-based byte-level LLMs can surpass subword-based LLMs of similar size, these models demand significantly greater computational resources and require considerably more time to train. In a related vein, \cite{el2020characterbert} introduced CharFormer, a BERT variant that constructs word representations via convolutions over character-level embeddings. While achieving notable gains in medical domain applications, CharFormer similarly incurs a higher computational cost. \cite{clark2022canine} propose CANINE, an encoder-only model with 150M parameters that operates at the character level. Central to CANINE is a deep transformer architecture—akin to our global model—augmented by a local transformer and strided convolutions for input downsampling. CANINE surpasses its token-level counterpart (mBERT) on several multilingual tasks. The ByT5 model~\citep{xue2022byt5} investigates encoder-decoder architectures at the byte level that eschew patching operations. While ByT5 demonstrates increased noise robustness and matches tokenizer-based models using only a quarter of the data, the absence of patching requires attention computations across all bytes in the sequence, substantially raising the computational burden. Modeling input sequences at the byte level increases their length, presenting obstacles for efficient scaling of such models. More recently, the Mamba Architecture~\citep{gu2023mamba}, which maintains a fixed-size memory state across extended contexts, was employed in \cite{wang2024mambabyte} to train byte-level Mamba models without patching. These models outperform byte-level transformers in flop-constrained environments at the 350M parameter scale, as measured by bits-per-byte performance across multiple datasets.

\textbf{Patching-based approaches:} By leveraging patching techniques, it is possible to mitigate the significant computational cost typically associated with byte-level LLMs, while still maintaining strong performance. Several studies have highlighted early positive outcomes in this area for relatively modest model sizes and limited training data. For example, \cite{nawrot-etal-2022-hierarchical} investigate a method involving fixed patching for both downsampling and upsampling, culminating in the hourglass transformer model, which surpasses other byte-level baselines when evaluated at the 150M parameter scale. Building upon this foundation, \cite{nawrot-etal-2023-efficient} propose enhancements through dynamic patching approaches. These involve a boundary predictor that can be learned in an end-to-end manner, a version guided by supervision from specific tokenizers, and an entropy-driven patching strategy akin to {\textsc BLT}; collectively, these refinements yield better results than standard transformers for language modeling at the 40M parameter level over datasets of about 400M tokens. In another line of work, \cite{lester2024training} explore the training of models on sequences compressed via arithmetic coding, allowing compression ratios to exceed those provided by BPE. Employing an equal-info windows method, their findings indicate improvements over byte-level baselines in language modeling, though these results still lag behind those of subword-based baselines.

Our research is heavily influenced by and most similar to MegaByte~\citep{yu2023megabyte}, a decoder-only causal LLM that employs a predetermined static approach for dividing and concatenating byte representations into patches, followed by a local model on the decoder end that reconstructs patches back into bytes. MegaByte is shown to achieve performance comparable to that of tokenizer-based models at the 1B parameter level using a dataset of approximately 400B bytes. In our experiments, we comprehensively examine MegaByte, noting that its static patching approach does not yet attain the efficiency of today’s best compute-optimized tokenizer-based models under fixed FLOPs constraints. Our work demonstrates that {\textsc BLT} addresses and narrows this performance gap. Similarly, \cite{slagle2024spacebyte} observe the same performance limitation in MegaByte, proposing an extension of the static patching method to operate on whitespaces and analogous space-like bytes, along with the addition of a local encoder model. Their results show performance gains over tokenizer-based transformer models within compute-equivalent conditions for certain domains, such as Github and arXiv, at the 1B parameter level. We also include empirical results with this model in our study, concluding that further architectural advancements are crucial for scaling byte-level models to achieve parity with the latest high-performing token-based models like Llama 3.

\section{Limitations and Future Work}

For this study, we adopt the optimal training step counts as established for Llama~3~\citep{dubey2024llama} to inform our architectural decisions. It should be noted, however, that these scaling laws were derived specifically for BPE-level transformer architectures and may not yield ideal data-to-parameter size ratios when applied to {\textsc BLT}. Determining the scaling laws tailored for {\textsc BLT} and assessing whether they confer more advantageous scaling behaviors is reserved for subsequent investigation. Furthermore, since the bulk of our experiments have been limited to model scales up to 1B parameters, the optimal configuration choices may shift with expansion to 8B parameters or greater. Such scaling may, in turn, enable superior performance for larger models.

Current transformer libraries and codebases are primarily optimized for tokenizer-based transformer models, achieving high efficiency within that context. Although our study provides theoretical flop-matched experiments and leverages efficient methods like FlexAttention for accommodating layers that differ from the standard transformer structure, our implementations might not fully match tokenizer-based models regarding wall-clock performance and could still realize improvements through additional optimization efforts.

Although {\textsc BLT} currently employs a separately pre-trained entropy model for patching, an intriguing avenue for future exploration lies in jointly training the patching mechanism in an end-to-end manner. Section \ref{sec:distillation} discusses preliminary studies that demonstrate promising outcomes for converting large-scale tokenizer-based architectures like Llama 3—which are pretrained on datasets exceeding 10T tokens—into byte-based models. This is achieved by initializing and subsequently freezing the global transformer using the original model weights. Continued investigation may yield strategies that not only preserve the advantages gained through bytefying, but also achieve superior performance relative to these original tokenizer-based models, without requiring training from the ground up.

\section{Conclusion}
\label{section:conclusion}

In this work, we introduce the Byte Latent Transformer (\textbf{BLT}), an innovative model architecture that challenges the prevailing reliance on fixed-vocabulary tokenization in large-scale language models. Through the implementation of a trainable, adaptive mechanism that organizes bytes into patches, {\textsc BLT} optimizes computational allocation according to the intrinsic complexity of the input data, thereby enhancing both computational efficiency and resilience. Our comprehensive scaling experiments reveal that {\textsc BLT} attains comparable performance to established tokenization-driven models such as Llama 3, even at large scales of up to 8B parameters and 4T processed bytes, and is able to decrease inference flops by as much as 50\% with only minor sacrifices in key evaluation metrics. In addition, {\textsc BLT} facilitates the simultaneous enlargement of model and patch sizes while maintaining constant inference costs, thereby opening up novel avenues for scaling, particularly well-suited to real-world computational scenarios. By directly modeling unprocessed byte sequences, {\textsc BLT} also enhances the system's robustness against input noise and grants improved sensitivity to rare linguistic phenomena, including richer modeling of sub-word information. Taken together, these findings establish {\textsc BLT} as a compelling, scalable, and resilient alternative to existing tokenization-centric models, presenting an adaptable and efficient foundation for next-generation language modeling.

\end{document}